\documentclass[11pt,a4paper]{article}

\usepackage[margin=2.6cm]{geometry}
\usepackage{lmodern}
\usepackage[T1]{fontenc}
\usepackage[utf8]{inputenc}
\usepackage{microtype}
\usepackage{xcolor}
\usepackage{booktabs}
\usepackage{listings}
\usepackage{enumitem}
\usepackage[colorlinks=true,linkcolor=calangoteal,urlcolor=calangoteal]{hyperref}
\usepackage{parskip}

\definecolor{calangoteal}{RGB}{0,110,110}
\definecolor{codebg}{RGB}{245,246,248}
\definecolor{codeframe}{RGB}{210,214,220}
\definecolor{shadegray}{RGB}{90,96,104}

\lstdefinestyle{shell}{
  basicstyle=\ttfamily\small,
  backgroundcolor=\color{codebg},
  frame=single,
  rulecolor=\color{codeframe},
  framesep=6pt,
  xleftmargin=8pt,
  xrightmargin=8pt,
  breaklines=true,
  breakatwhitespace=true,
  postbreak=\mbox{\textcolor{shadegray}{$\hookrightarrow$}\space},
  columns=fullflexible,
  keepspaces=true,
  showstringspaces=false,
  upquote=true,
}
\lstset{style=shell}

\newcommand{\file}[1]{\texttt{#1}}
\newcommand{\opt}[1]{\texttt{#1}}

\title{\bfseries Calango Packaging Guide\\[4pt]
\Large Building the macOS \file{.dmg} and Linux \file{.deb} installers}
\author{Seixas Research}
\date{July 2026 --- Calango 26.7.13}

\begin{document}
\maketitle

\begin{abstract}
\noindent
Calango ships as two native installers built with CMake/CPack: a
drag-and-drop disk image (\file{.dmg}) for macOS and a Debian package
(\file{.deb}) for Ubuntu/Debian. This guide covers the prerequisites, the
exact commands, what the packaging pipeline does under the hood, how to
verify the artifacts, and the failure modes we have already hit --- with
their fixes. Everything described here lives in \file{CMakeLists.txt}
(install and CPack sections) and in the \file{packaging/} directory of the
repository.
\end{abstract}

\tableofcontents
\bigskip
\hrule

\section{Overview}

Both installers come out of the standard CMake flow: \emph{configure}
$\to$ \emph{build} $\to$ \emph{cpack}. The platform decides the generator:

\begin{center}
\small
\begin{tabular}{@{}lll@{}}
\toprule
Platform & CPack generator & Artifact \\
\midrule
macOS (\opt{-DCALANGO\_MACOS\_BUNDLE=ON}) & \opt{DragNDrop} &
\file{calango-26.7.13-macos-arm64.dmg} \\
Linux & \opt{DEB} & \file{calango\_26.7.13\_amd64.deb} \\
\bottomrule
\end{tabular}
\end{center}

Each artifact is accompanied by a \file{.sha256} checksum file. The
package version comes from \opt{project(calango VERSION ...)} in
\file{CMakeLists.txt} --- the single source of truth for the version,
baked into the binary as \opt{CALANGO\_VERSION} (About dialog, project
files) and reused by CPack for the installer file names.

The supporting files are:

\begin{itemize}[itemsep=2pt]
  \item \file{packaging/macos/make\_icns.sh} --- renders
        \file{calango.icns} from the master PNG icon in
        \file{assets/calango/} (used when no committed \file{.icns}
        exists).
  \item \file{packaging/macos/embed\_python\_framework.sh} --- makes the
        \file{.app} independent of the build machine's Python
        (Section~\ref{sec:python-embed}).
  \item \file{packaging/linux/calango.desktop} --- application launcher entry.
  \item \file{packaging/linux/calango-mime.xml} --- registers the
        \texttt{application/x-calango-project} MIME type for
        \file{.calproj} project files.
  \item \file{packaging/linux/postinst}, \file{postrm} --- refresh the
        desktop/MIME/icon caches on install and removal.
\end{itemize}

\section{macOS: building the \file{.dmg}}

\subsection{Prerequisites}

\begin{itemize}[itemsep=2pt]
  \item Xcode command-line tools (provides \file{sips}, \file{iconutil},
        \file{codesign}, \file{hdiutil}).
  \item Qt 6 with the Svg module; Homebrew's \lstinline|brew install qt|
        puts it at \file{/opt/homebrew/opt/qt}.
  \item CMake $\geq$ 3.21.
  \item A Python virtualenv with ASE installed, e.g.\ the project
        \file{.venv} (\lstinline|python3 -m venv .venv && .venv/bin/pip install ase|).
\end{itemize}

\subsection{Commands}

From the repository root:

\begin{lstlisting}
cmake -S . -B build-dmg \
  -DCMAKE_BUILD_TYPE=Release \
  -DPython3_EXECUTABLE=$PWD/.venv/bin/python \
  -DCMAKE_PREFIX_PATH=/opt/homebrew/opt/qt \
  -DCALANGO_MACOS_BUNDLE=ON

cmake --build build-dmg -j 8

cd build-dmg && cpack
\end{lstlisting}

The result is \file{build-dmg/calango-26.7.13-macos-arm64.dmg}. Mounting
it shows \file{calango.app} next to an \file{/Applications} symlink ---
the standard drag-to-install layout. There is deliberately no
click-through license sheet
(\opt{CPACK\_DMG\_SLA\_USE\_RESOURCE\_FILE\_LICENSE OFF}): a license SLA
blocks scripted \file{hdiutil attach} and adds friction for an
MIT-licensed application.

\subsection{What the install step does}

\opt{cpack} internally runs the CMake install rules into a staging
directory (\file{\_CPack\_Packages/}) and then wraps the result in a
disk image. The install rules perform, in order:

\begin{enumerate}[itemsep=2pt]
  \item \textbf{Stage the bundle.} \file{calango.app} is installed with
        the Finder icon (\file{calango.icns}, committed or rendered at
        build time by \file{make\_icns.sh}) and the \file{Info.plist}
        metadata (bundle id \texttt{org.seixasresearch.calango}).
  \item \textbf{Optionally embed a Python environment} into
        \file{Contents/Resources/python}, when the cache variable
        \opt{CALANGO\_EMBEDDED\_PYTHON\_DIR} is set
        (Section~\ref{sec:python-embed}).
  \item \textbf{Deploy Qt with \file{macdeployqt} --- twice.} The tool
        copies the Qt frameworks and plugins into the bundle and rewrites
        their load paths to \texttt{@executable\_path}. Two passes with
        \opt{-libpath=/opt/homebrew/opt/qt/lib} are required: frameworks
        copied into the bundle lose the \texttt{@loader\_path} context
        their \texttt{@rpath} dependencies resolve against, so
        transitive dependencies (e.g.\ QtGui $\to$ QtDBus, plugin
        dependencies) are only found when the second pass re-scans the
        already-copied frameworks. The first pass prints a wall of
        \texttt{ERROR: Cannot resolve rpath} lines --- this is expected
        noise; judge success by the verification steps below.
  \item \textbf{Embed the linked \file{Python.framework}.}
        \file{embed\_python\_framework.sh} copies the framework version
        the binary links against (typically Homebrew's) into
        \file{Contents/Frameworks}, replaces Homebrew's dangling
        \file{site-packages} symlink with a real directory, recreates the
        canonical framework symlinks, retargets the load command with
        \file{install\_name\_tool}, and finally re-signs the whole bundle
        ad hoc (\lstinline|codesign --force --deep --sign -|). Without
        the re-sign, Apple-silicon macOS refuses to load the modified
        binaries.
\end{enumerate}

\subsection{Verification}

\begin{lstlisting}
APP=_CPack_Packages/Darwin/DragNDrop/calango-26.7.13-macos-arm64/calango.app

# 1. Signature must verify strictly, including nested code
codesign --verify --deep --strict "$APP"

# 2. The packaged binary must start and find a Python with ASE
"$APP/Contents/MacOS/calango" --probe-python

# 3. The image must mount without an SLA prompt
hdiutil attach -readonly -nobrowse calango-26.7.13-macos-arm64.dmg
\end{lstlisting}

To rebuild the package from a clean state, clear the staging area first
with \lstinline|cmake -E rm -rf _CPack_Packages| --- stale staging can
mask deployment bugs.

\subsection{Troubleshooting}

\begin{center}
\small
\begin{tabular}{@{}p{0.44\textwidth}p{0.5\textwidth}@{}}
\toprule
Symptom & Cause / fix \\
\midrule
\texttt{ERROR: Cannot resolve rpath "@rpath/Qt...''} during
\opt{cpack} &
Expected noise from \file{macdeployqt}'s first pass. Only investigate if
verification fails afterwards. \\
\addlinespace
\texttt{dyld: Library not loaded: @rpath/QtDBus...} when launching the
packaged app &
A framework is missing from the bundle --- the second
\file{macdeployqt} pass or the \opt{-libpath} argument was removed, or
the install rpath was stripped. Both fixes live in the
\opt{install(CODE ...)} block of \file{CMakeLists.txt}. \\
\addlinespace
\texttt{codesign} reports \texttt{No such file or directory} for a path
that exists &
A dangling symlink inside the bundle (Homebrew's
\file{site-packages} link is the known case, handled by
\file{embed\_python\_framework.sh}). Find offenders with
\lstinline|find app -type l ! -exec test -e {} \; -print|. \\
\addlinespace
\texttt{hdiutil: attach canceled} &
The DMG carries a click-through SLA. Keep
\opt{CPACK\_DMG\_SLA\_USE\_RESOURCE\_FILE\_LICENSE} set to \opt{OFF}. \\
\addlinespace
\texttt{Permission denied} running a \file{packaging/} script &
Cloud-synced checkouts (Dropbox) strip executable bits. CMake therefore
invokes every packaging script via \lstinline|bash <script>|; keep that
pattern when adding new ones. \\
\bottomrule
\end{tabular}
\end{center}

\section{Linux: building the \file{.deb}}

\subsection{Prerequisites (Ubuntu/Debian)}

\begin{lstlisting}
sudo apt install build-essential cmake ninja-build \
  qt6-base-dev libqt6svg6-dev libqt6opengl6-dev \
  libgl1-mesa-dev python3-dev python3-venv \
  dpkg-dev file

python3 -m venv .venv && .venv/bin/pip install ase
\end{lstlisting}

\opt{dpkg-dev} provides \file{dpkg-shlibdeps}, which CPack uses to
compute the runtime \texttt{Depends:} line automatically
(\opt{CPACK\_DEBIAN\_PACKAGE\_SHLIBDEPS ON}).

\subsection{Commands}

\begin{lstlisting}
cmake -S . -B build-deb \
  -DCMAKE_BUILD_TYPE=Release \
  -DPython3_EXECUTABLE=$PWD/.venv/bin/python

cmake --build build-deb -j$(nproc)

cd build-deb && cpack
\end{lstlisting}

The result is \file{calango\_26.7.13\_<arch>.deb}. Install it via
\opt{apt} rather than \opt{dpkg -i}, so the dependencies are resolved:

\begin{lstlisting}
sudo apt install ./calango_26.7.13_amd64.deb
\end{lstlisting}

\subsection{Package contents and desktop integration}

\begin{center}
\small
\begin{tabular}{@{}lp{0.32\textwidth}@{}}
\toprule
Path in package & Purpose \\
\midrule
\file{/usr/bin/calango} & The application binary \\
\file{/usr/share/applications/calango.desktop} & Launcher entry
(\texttt{Exec=calango \%F}) \\
\file{/usr/share/mime/packages/calango-mime.xml} & MIME type for
\file{.calproj} \\
\file{/usr/share/icons/hicolor/scalable/apps/calango.svg} & Scalable icon \\
\file{/usr/share/pixmaps/calango.png} & Legacy raster icon \\
\file{/usr/lib/calango/python/} & Embedded Python (optional, see
Section~\ref{sec:python-embed}) \\
\bottomrule
\end{tabular}
\end{center}

The MIME registration means double-clicking a \file{.calproj} workspace
in the file manager opens it in Calango; \file{MainWindow::loadFile}
routes \file{.calproj} arguments to the project loader rather than the
ASE structure reader. The \file{postinst}/\file{postrm} maintainer
scripts refresh \file{update-mime-database},
\file{update-desktop-database} and \file{gtk-update-icon-cache} so the
association is live immediately after install (modern dpkg triggers
usually do this anyway; the scripts make it robust on minimal systems).

Unless an embedded Python is bundled, the package declares
\texttt{Depends: python3 (>= 3.9)} and \texttt{Recommends: python3-ase}.
Without ASE the GUI starts but structure I/O and job submission are
disabled; users can also point Calango at any interpreter via
\texttt{CALANGO\_PYTHON}.

\subsection{Verification}

\begin{lstlisting}
dpkg-deb --info  calango_26.7.13_amd64.deb   # control fields, Depends
dpkg-deb --contents calango_26.7.13_amd64.deb

# after installing:
calango --probe-python
xdg-mime query default application/x-calango-project
\end{lstlisting}

\noindent
\textbf{Status note:} the DEB configuration is complete but has so far
been validated only structurally (the development machine is a Mac
without \file{dpkg}). The first CI or VM build on Ubuntu should run the
verification steps above.

\section{The embedded Python environment}
\label{sec:python-embed}

Calango embeds CPython at runtime. The interpreter is resolved in this
order (documented in \file{src/python\_bridge/PythonEngine.hpp}):

\begin{enumerate}[itemsep=1pt]
  \item \texttt{\$CALANGO\_PYTHON} --- explicit interpreter path;
  \item \texttt{\$VIRTUAL\_ENV/bin/python} --- active virtualenv;
  \item a Python bundled by the installers:
        \file{Contents/Resources/python/bin/python3} (macOS) or
        \file{<prefix>/lib/calango/python/bin/python3} (Linux);
  \item the interpreter CMake was configured against
        (\opt{CALANGO\_DEFAULT\_PYTHON}, meaningful only on the build
        machine).
\end{enumerate}

To ship a fully self-contained installer, pass a \emph{relocatable}
Python distribution that has ASE installed:

\begin{lstlisting}
cmake -S . -B build-dmg ... \
  -DCALANGO_EMBEDDED_PYTHON_DIR=/path/to/standalone-python
\end{lstlisting}

A distribution from the \texttt{python-build-standalone} project is the
recommended source. A plain project \file{.venv} is \emph{not}
relocatable --- its \file{bin/python} and \file{pyvenv.cfg} refer back to
the build machine's base interpreter.

Independent of this option, the macOS pipeline always embeds the
\file{Python.framework} that the binary itself links against
(\file{libpython} via \opt{Development.Embed}), so the packaged app
launches on machines without Homebrew. For the interpreter to also find
ASE there, either bundle an environment as above or let users install
ASE and set \texttt{CALANGO\_PYTHON}.

\section{Known limitations and next steps}

\begin{itemize}[itemsep=2pt]
  \item \textbf{Code signing:} the \file{.app} is signed ad hoc.
        Distribution outside a lab requires a Developer~ID certificate
        and notarization (\file{codesign --sign "Developer ID..."},
        \file{notarytool submit}, \file{stapler}); the ad-hoc step in
        \file{embed\_python\_framework.sh} is the place to swap the
        identity in.
  \item \textbf{DEB validation:} run the Linux pipeline once in CI
        (Ubuntu 22.04/24.04) and consider \file{lintian} for policy
        checks.
  \item \textbf{Full ASE bundling:} wiring a
        \texttt{python-build-standalone} + ASE tree into both installers
        (and building against its \file{libpython}) would remove the
        last runtime dependency on a user-provided Python.
\end{itemize}

\end{document}
